\subsection{Development and reasoning exercises}

\begin{enumerate}
\item Draw one point in both Cartesian and polar coordinates. Derive
\[
x=r\cos\theta,
\qquad
y=r\sin\theta,
\qquad
r=\sqrt{x^2+y^2},
\]
and explain why the polar description is better adapted to a distinguished focus.

\item Starting from
\[
d(P,F)=e\,d(P,D),
\]
place the pole at $F$ and the directrix at $x=d$. Derive
\[
r=\frac{p}{1+e\cos\theta},
\qquad p=ed,
\]
using a fully labelled geometric diagram.

\item Fix $p=2$ and compare the ellipses obtained for $e=0.2$, $e=0.6$, and $e=0.9$. Explain the changes using $r_{\min}$ and $r_{\max}$ rather than only visual appearance.

\item Explain directly from the denominator why $e=1$ is the transition between a closed ellipse and an open hyperbola.

\item The polar formula with $e>1$ and $r\geq0$ describes one hyperbolic branch relative to one focus and directrix. Explain how the complete two-branch hyperbola is recovered and why the Cartesian equation displays it more naturally.

\item Compare
\[
\frac{x^2}{a^2}+\frac{y^2}{b^2}=1
\]
with
\[
r=\frac{a(1-e^2)}{1+e\cos\theta}.
\]
State which geometric information is immediate in each description.

\item Explain why the parameter $t$ in
\[
x=a\cos t,
\qquad y=b\sin t
\]
is not generally the same as the polar angle $\theta$ measured from a focus.

\item An orbit has nearest focal distance $4$ and farthest focal distance $12$. Find $a$, $e$, and $p$, then write its polar equation with periapsis at $\theta=0$.

\item Explain how $p$, $e$, and $\theta_0$ separately control scale, shape, and orientation in
\[
r=\frac{p}{1+e\cos(\theta-\theta_0)}.
\]

\item Summarise the chapter as a change of viewpoint rather than a classification table: Cartesian geometry first, polar coordinates second, one focus--directrix law last. Explain why this sequence prepares a later derivation of orbits in mechanics.
\end{enumerate}
